\documentclass[11pt]{article}
\usepackage{preamble}
\setlength{\parskip}{4pt}

\begin{document}

\daybanner{AP Statistics \textperiodcentered\ Unit 1 \textperiodcentered\ Topic 1.10 \textperiodcentered\ B: 2026-09-28 \textperiodcentered\ E: 2026-10-05 \textperiodcentered\ TEACHER KEY}{What Can Each Study Claim?}

\sectionbanner{CED TETHER}

{\small
\begin{itemize}[leftmargin=1.5em,itemsep=2pt,topsep=2pt]
  \item \textbf{Skill 1.C} --- Describe an appropriate method for gathering and representing data.
  \item \textbf{Skill 4.A} --- Make an appropriate claim or draw an appropriate conclusion.
  \item \textbf{EU DAT-2} --- The way we collect data influences what we can and cannot say about a population.
  \item \textbf{LO DAT-2.A} --- Identify the type of a study. [Skill 1.C]
  \item \textbf{EK DAT-2.A.1} --- A population consists of all items or subjects of interest.
  \item \textbf{EK DAT-2.A.2} --- A sample selected for study is a subset of the population.
  \item \textbf{EK DAT-2.A.3} --- In an observational study, treatments are not imposed. Investigators examine data for a sample of individuals (retrospective) or follow a sample of individuals into the future collecting data (prospective) in order to investigate a topic of interest about the population. A sample survey is a type of observational study that collects data from a sample in an attempt to learn about the population from which the sample was taken.
  \item \textbf{EK DAT-2.A.4} --- In an experiment, different conditions (treatments) are assigned to experimental units (participants or subjects).
  \item \textbf{LO DAT-2.B} --- Identify appropriate generalizations and determinations based on observational studies. [Skill 4.A]
  \item \textbf{EK DAT-2.B.1} --- It is only appropriate to make generalizations about a population based on samples that are randomly selected or otherwise representative of that population.
  \item \textbf{EK DAT-2.B.2} --- A sample is only generalizable to the population from which the sample was selected.
  \item \textbf{EK DAT-2.B.3} --- It is not possible to determine causal relationships between variables using data collected in an observational study.
\end{itemize}
}
\textbf{Essential question:} How do selection and assignment determine the conclusions a study can support?\par
\textbf{DOK-3 skill:} 4.A \textperiodcentered\ \textbf{FRQ pattern:} {\small\texttt{observational-\allowbreak{}vs-\allowbreak{}experiment-\allowbreak{}what-\allowbreak{}each-\allowbreak{}can-\allowbreak{}claim}}\par
\textbf{DOK rationale:} Part (c) weighs two study plans against a causal claim, connects self-selection to a concrete lurking variable, and separately bounds what random assignment among volunteers can establish and generalize.\par

\frameworkphaseheader{Do Now (first take) $\rightarrow$ Explore (video + follow-along) $\rightarrow$ Exit (finish + turn in)}{1 $\rightarrow$ 3}{45}{%
\begin{itemize}[leftmargin=1.5em,itemsep=2pt,topsep=2pt]
  \item Board slide up at the bell. Ask for a plan choice before defining study types.
  \item Begin the video follow-along at minute 5.
  \item Return to the board slide at minute 33. Circulate and separate assignment from selection.
  \item Collect at the door. Score part (c) only, E/P/I.
\end{itemize}
}{%
\begin{itemize}[leftmargin=1.5em,itemsep=2pt,topsep=2pt]
  \item Choose Plan A or Plan B and name one design feature in a first take.
  \item Complete the video follow-along about populations, samples, observational studies, and experiments.
  \item Finish (a)--(c), revise the first take in the exit line, and turn in the sheet.
\end{itemize}
}{%
\begin{itemize}[leftmargin=1.5em,itemsep=2pt,topsep=2pt]
  \item Who decides whether a student receives tutoring in each plan?
  \item How could motivation affect both the group a student joins and the score later measured?
  \item What does the coin flip balance, and among whom?
  \item Were the volunteers selected by chance from all students? What conclusion does that limit?
\end{itemize}
}{Solo teacher. Circulate ~5 pairs. Ask students to label each arrow as selection or assignment before accepting a claim about cause or scope.}

\begin{callout}{\IconWarn\ WATCH FOR}{calloutred}
\begin{itemize}[leftmargin=1.5em,itemsep=2pt,topsep=2pt]
  \item Calling any comparison of two groups an experiment.
  \item Naming a lurking variable without linking it to both explanatory and response variables.
  \item Assuming random assignment also makes volunteers representative of all students.
\end{itemize}
\end{callout}

\sectionbanner{SCORING PART (c) --- E / P / I}

\textbf{Expected elements}
\begin{itemize}[leftmargin=1.5em,itemsep=2pt,topsep=2pt]
  \item \textbf{random-assignment} (required) --- Chooses B and links random assignment to balancing outside factors
  \item \textbf{alternative-cause} (required) --- Connects a plausible outside factor to both tutoring choice and scores in A
  \item \textbf{volunteer-scope} (required) --- Explains why assigning the volunteers does not make them representative of all students
\end{itemize}

{\small\renewcommand{\arraystretch}{1.25}
\noindent\begin{tabular}{|p{0.5in}|p{5.9in}|}
\hline
\textbf{E} & All three elements are correct and linked to this problem. Accept equivalent plain-language reasoning. \\ \hline
\textbf{P} & At least one element includes a correct evidence-to-reason link, but another is missing or incorrect. \\ \hline
\textbf{I} & No correct evidence-to-reason link; only a choice, copied numbers, or unsupported statements. \\ \hline
\end{tabular}}

\textbf{Common mistakes}
\begin{itemize}[leftmargin=1.5em,itemsep=2pt,topsep=2pt]
  \item Calling Plan A an experiment because it compares two groups
  \item Treating the coin flip as random selection from the school instead of random assignment among volunteers
  \item Naming motivation without explaining how it relates to both attendance and scores
  \item Generalizing Plan B to all students solely because it is an experiment
\end{itemize}

\clearpage
\focusbanner{TODAY'S PROBLEM --- annotated}

Suppose a student wants to know whether a new tutoring center raises test scores.\par\smallskip \textbf{Plan A:} Compare this year's test scores for students who choose to attend the center with scores for students who do not attend.\par \textbf{Plan B:} Recruit 40 student volunteers, then use a coin flip to assign each volunteer either to attend the center or not to attend. Compare their later test scores.\par
\begin{center}
\textbf{Two plans for studying tutoring}\par\smallskip
\begin{tabular}{|l|c|c|c|}
\hline
\textbf{Plan} & \textbf{Who?} & \textbf{Groups?} & \textbf{Response} \\ \hline
\textbf{A} & Students & Choose & Scores \\ \hline
\textbf{B} & 40 volunteers & Coin flip & Scores \\ \hline
\end{tabular}
\end{center}
\begin{firsttakebox}{FIRST TAKE (5 min)}
Which plan gives stronger evidence that the center raises scores? Name the feature behind your choice.\par
\answer{Not graded. Listen for students who say Plan B is stronger merely because it is random without naming what is randomized.}
\end{firsttakebox}

\textit{Rules callout ``SELECTION, ASSIGNMENT, AND SCOPE (keep on the page)'' is printed on the student sheet.}\par\medskip

\dokbadge{1}~\textbf{(a)} Identify each plan as an observational study or an experiment.\par
\answer{Plan A is observational because students choose whether to attend; no treatment is assigned. Plan B is an experiment because attendance is assigned to volunteers by a coin flip.}

\dokbadge{2}~\textbf{(b)} Name one outside factor that could affect both choosing tutoring and test scores in Plan A. Explain both links in two sentences.\par
\answer{One example is motivation: more motivated students may be more likely to choose tutoring and may also study more or perform better even without the center. Prior achievement, available time, or family support are acceptable if both links are explained.}

\dokbadge{3}~\textbf{(c)} Which plan offers stronger evidence that tutoring causes a change in scores? Explain why its design helps with the outside factor from (b). Then state why its result may not apply to all students at the school.\par
\answer{Plan B could support a causal claim for these volunteers because random assignment tends to balance lurking variables between the attendance and no-attendance groups. Under Plan A, a score difference could be due to motivation or another preexisting difference rather than tutoring. Because the 40 participants volunteered rather than being randomly selected from a named population, Plan B cannot automatically generalize the effect to all students at the school or to students elsewhere.}

\begin{summaryexitbox}{EXIT}
One thing the video changed about my first take: \hrulefill
\end{summaryexitbox}

\end{document}
